\section{Parton distributions}

\subsection{Generic matrix element and
Lorentz invariance}

Historically \cite{Feynman:1973xc}, PDFs were introduced
to describe \mbox{spin-1/2 quarks.}
Since complications related to spin
do not affect
the very concept of parton distributions, we start with
a simple example of a scalar theory.
In that case, information about the target is accumulated
in the generic matrix element $\langle p | \phi (0) \phi(z) | p \rangle $.
By Lorentz invariance, it is a function of two scalars, $(pz) \equiv - \nu$ and $z^2$
(or $-z^2$ if we want a positive value for spacelike $z$):
\begin{align}
\langle p |   \phi(0) \phi (z)|p \rangle
=  & {\cal M} (-(pz), -z^2)
\, .
\label{lorentz}
\end{align}
It can be shown \cite{Radyushkin:2016hsy,Radyushkin:1983wh} that, for all contributing Feynman diagrams,
its Fourier transform ${\cal P} (x, -z^2)$ with respect to $(pz)$
has the $-1 \leq x \leq 1$ support, i.e.,
\begin{align}
{\cal M} (-(pz), -z^2)
&   =
\int_{-1}^1 dx
\, e^{-i x (pz) } \,  {\cal P} (x, -z^2)  \   .
\label{MPD}
\end{align}
Note that Eq. (\ref{MPD}) gives a covariant definition of $x$.
There is no need to assume that
$p^2=0$ or $z^2=0$ to define $x$.

\subsection{Collinear PDFs}

Choosing
some special cases of $p$ and $z$, one can get expressions for various parton distributions,
all in terms of the same function ${\cal M} (-(pz), -z^2)$.
In particular, taking a light-like $z$, e.g., that
having the light-front minus component $z_-$ only, we
parameterize the matrix element by the twist-2 parton distribution $f(x)$
\begin{align}
{\cal M} (-p_+ z_- , 0)  =
\int_{-1}^1 dx \, f(x) \,
e^{-ixp_+ z_-} \,  \   ,
\label{twist2par0}
\end{align}
with $f(x)$ having the usual interpretation of probability that
the parton carries the fraction $x$ of the target momentum
component $p_+$.
The inverse relation is given by
\begin{align}
f  (x)  =\frac{1}{2 \pi}  \int_{-\infty}^{\infty}  d\nu \, e^{-i x \nu}  \, {\cal M}(\nu , 0)    = {\cal P} (x, 0)
\  .
\label{fxMnu}
\end{align}
Since $f(x) =   {\cal P} (x, 0) $, the function $ {\cal P} (x, -z^2) $ generalizes
PDFs onto non-lightlike intervals $z^2$, and we will call it
{\it pseudo-PDF}. The variable $(pz)$ is called
the {\it Ioffe time} \cite{Ioffe:1969kf},
and ${\cal M}(\nu , -z^2) $ is the {\it Ioffe-time distribution} \cite{Braun:1994jq}.

Note that the definition of $ {\cal P} (x, -z^2) $ is simpler than that of $f(x)$ because it
does not require taking a subtle $z^2 \to 0$ limit.
In renormalizable theories,
the function $ {\cal M} (\nu, z^2)  $
has $\sim \ln z^2$ singularities
generating perturbative evolution of parton densities.
Within the operator product expansion
(OPE) approach, the $\ln z^2$ singularities
are subtracted using some prescription, say, dimensional renormalization,
and the resulting PDFs depend on the renormalization scale $\mu$, i.e.,
$f(x) \to   f(x, \mu^2)$.

\subsection{Transverse momentum dependent
distributions}

When $z^2$ is spacelike, one can treat $-z^2$
as the magnitude squared of a two-dimensional
vector $ \{z_1,z_2\}$, and introduce a two-dimensional
Fourier transform
with respect to its components, i.e., to write
\begin{align}
{\cal P} (x,  z_1^2+z_2^2)
&   =
\int_{-\infty}^\infty  {d k_1 } e^{i k_1 z_1}
\nn & \times
\int_{-\infty}^\infty  dk_2 \,e^{ik_2z_2}
{\cal F} (x, k_1^2+k_2^2)
\   .
\label{PTMD}
\end{align}
Due to rotational invariance of ${\cal P} (x,  z_1^2+z_2^2) $
in $ \{z_1,z_2\}$ plane,
the function ${\cal F} (x, k_1^2+k_2^2)$ depends
on ${k_1,k_2}$ through $k_1^2+k_2^2$, the fact already reflected in the notation.
Combining this representation with Eq. (\ref{MPD}), one has
\begin{align}
{\cal M} (\nu,  z_1^2+z_2^2)
&   =  \int_{-1}^1 dx
\, e^{i x \nu  } \,
\int_{-\infty}^\infty  {d k_1 } e^{i k_1 z_1}
\nn & \times
\int_{-\infty}^\infty  dk_2 \,e^{ik_2z_2}
{\cal F} (x, k_1^2+k_2^2)
\   .
\label{MTMD}
\end{align}

A physical interpretation of ${\cal F} (x, k_1^2+k_2^2)$
may be given in the frame where the target momentum $p$ is longitudinal,
\mbox{$p= (E, {\bf 0}_\perp, P)$,} while the vector
$\{z_1,z_2\} $ is in the transverse plane. Taking $z$ that has
$z_-$ and $z_\perp$ components only, one can identify ${\cal F} (x, k_\perp^2)$
with the
{\it TMD } and write
\begin{align}
{\cal P} (x,  z_\perp^2)
&   =
\int  d^2{\bf k}_\perp   e^{i ({\bf k}_\perp {\bf z}_\perp)}
{\cal F} (x, {k}_\perp^2)
\   .
\label{MTMD0}
\end{align}
In this case, the pseudo-PDFs $ {\cal P} (x,  z_\perp^2) $ coincide
with the {\it impact parameter distributions}, a well-known concept
actively
used in TMD studies.

The \mbox{$\sim \ln z_\perp^2$} terms
in ${\cal M} (\nu,  z_\perp^2) $ are
produced by the $\sim 1/k_\perp^2$ hard tail of ${\cal F} (x, k_\perp^2)$.
Thus, it makes sense to visualize
${\cal M} (\nu, z_\perp^2)  $ as a sum of a soft part ${\cal M}^{\rm soft} (\nu, z_\perp^2)$,
that has a finite $z_\perp^2\to 0$ limit
and a hard part reflecting the evolution.
For TMDs, soft part decreases faster than $1/k_\perp^2$, say, like
a Gaussian $e^{-k_\perp^2/\Lambda^2}$. In the $z_\perp$ space, the distributions
are then concentrated in \mbox{$z_\perp \sim 1/\Lambda$} region.
